\chapter{\texorpdfstring{A criterion for $\Om_n \cong \Ct{n}$}{A criterion for the full wreath
  power}}

% The counter is set so that the items below carry the numbers 1.0--1.6 of the paper.
\setcounter{theorem}{-1}

In this section, we derive criterion $(*)$ from the introduction. To begin with, we list a few
simple

\begin{facts}[Kummer tower]
  \label{facts:1.0}
  \uses{def:iterates, def:galois}
  \lean{QuadraticIterates.natDegree_iteratedPoly, QuadraticIterates.kummer_degree_bound,
    QuadraticIterates.splittingField_le_succ, QuadraticIterates.splittingField_mono,
    QuadraticIterates.relfinrank_succ_le,
    QuadraticIterates.relfinrank_succ_eq_finrank_adjoin,
    QuadraticIterates.mem_rootSet_iteratedPoly_succ, QuadraticIterates.exists_sq_eq_sub}
  \leanok
  $f_n$ has degree $2^n$. $K_{n+1} = K_n\bigl(\sqrt{\alpha - a} \bigm| \alpha \text{ root of }
  f_n\bigr)$, so $K_n \subset K_{n+1}$ is a $2$-Kummer extension, and
  $[K_{n+1} : K_n] \le 2^{2^n}$.
\end{facts}
\begin{proof}
  \leanok
  \uses{lem:multiquadratic}
  The degree statement is immediate from $f_{n+1} = f_n^2 + a$. A root of $f_{n+1}$ is precisely
  an element $\beta$ with $\beta^2 + a$ a root of $f_n$; as $\Qbar$ is algebraically closed,
  every root $\alpha$ of $f_n$ has such a $\beta$, namely a square root of $\alpha - a$. Hence
  $K_{n+1}$ is generated over $K_n$ by square roots of the $2^n$ elements $\alpha - a \in K_n$,
  and a multiquadratic extension generated by $r$ square roots has degree at most $2^r$
  (Lemma~\ref{lem:multiquadratic}).
\end{proof}

\begin{lemma}[cf.\ \cite{odoni}, Lemma 4.4]
  \label{lem:1.1}
  \uses{def:cseq, def:bseq}
  \lean{QuadraticIterates.cSeq_pos, QuadraticIterates.abs_le_cSeq,
    QuadraticIterates.cSeq_gcd, QuadraticIterates.cSeq_associated_gcd,
    QuadraticIterates.intCast_bSeq, QuadraticIterates.cSeq_eq_prod_bSeq,
    QuadraticIterates.isCoprime_bSeq, QuadraticIterates.bSeq_ne_zero,
    QuadraticIterates.cSeq_ne_zero, QuadraticIterates.cSeq_factorization_shape}
  \leanok
  \begin{enumerate}
    \item[a)] $\forall n \ge 2 : c_n > 0$.
    \item[b)] The $b_n$ $(n \ge 1)$ are integers such that
      $\forall n \ge 1 : c_n = \prod_{d \dvd n} b_d$, and the $b_n$ are pairwise coprime.
  \end{enumerate}
\end{lemma}
\begin{proof}
  \leanok
  \uses{lem:strongdvd, lem:moebiusfactor}
  a) $a \ne 0, -1$ since $-a$ is not a square. Therefore $c_2 = a^2 + a \ge |a| > 0$, and
  induction shows that $c_{n+1} = c_n^2 + a \ge a^2 + a \ge |a| > 0$.

  b) Since all $c_n \ne 0$, the rational number $\prod_{d \dvd n} c_d^{\mu(n/d)}$ is defined, and
  we have to show that it is already an integer. Let $p$ be a prime dividing at least one of the
  $c_n$. Let $m := \min\{n \ge 1 \mid p \dvd c_n\}$ and $e := v_p(c_m)$. Then we have
  $p \dvd c_n$ iff $p^e \dvd c_n$ iff $m \dvd n$. Indeed, with $x_n := f_n(0)$ $(= \pm c_n$ for
  $n \ge 1)$ we have $x_m \equiv 0 \bmod p$ and $\bmod\ p^e$, but $\not\equiv 0 \bmod p^{e+1}$;
  therefore $(x_n)_{n \ge 0}$ is purely periodic $\bmod\ p$ and $\bmod\ p^{e+1}$ with period $m$,
  and mod $p^{e+1}$ we have $x_{m+1} = x_m^2 + a \equiv a \equiv x_1$, so $(x_n)$ is periodic
  from $n = 1$ on with period $m$, hence $x_n$ is never zero mod $p^{e+1}$ if $m \nmid n$. That
  is, $v_p(c_n) = e$ if $m \dvd n$, and $0$ otherwise. Now
  $v_p(b_n) = \sum_{d \dvd n} v_p(c_d)\,\mu(n/d) = e \sum_{d \dvd n/m} \mu(n/(dm)) = e$ if
  $n = m$, and $0$ otherwise. Therefore $b_n$ is integral at $p$, and moreover $p$ divides only
  one of the $b_n$. Hence all the $b_n$ are integers and pairwise coprime; Möbius inversion gives
  $c_n = \prod_{d \dvd n} b_d$.
  \formalnote{The formalization splits this argument into its two halves and proves each in the
  generality it deserves. Integrality (and Möbius inversion) come from \emph{strong
  divisibility}: $\gcd(c_m, c_n) = |c_{\gcd(m,n)}|$ (Lemma~\ref{lem:strongdvd}), which for a
  nowhere-vanishing sequence forces every Möbius quotient to be integral
  (Lemma~\ref{lem:moebiusfactor}) --- this replaces the periodicity computation above.
  Coprimality does use the valuation shape, but in the form ``$v_p(c_n) = E \cdot [m \dvd n]$ for
  some $m \ge 1$, $E$'' (\lname{cSeq\_factorization\_shape}), which is proved over an arbitrary
  UFD (Theorem~\ref{thm:valshape}) and fed to the general coprimality lemma
  \lname{moebiusFactorR\_isRelPrime}.}
\end{proof}

The next lemma gives us a simple irreducibility test for the $f_n$.

\begin{lemma}
  \label{lem:1.2}
  \uses{def:cseq, def:iterates}
  \lean{QuadraticIterates.irreducible_iteratedPoly_of_not_isSquare_cSeq,
    QuadraticIterates.even_reducible_factorization,
    QuadraticIterates.isUnit_or_associated_of_dvd_comp,
    QuadraticIterates.isUnit_or_associated_of_dvd_comp_of_associated,
    QuadraticIterates.map_iteratedPoly_comp_neg_X, QuadraticIterates.map_iteratedPoly_succ}
  \leanok
  If none of $c_1, c_2, \dots, c_n$ is a square in $\QQ$, then $f_n$ is irreducible in $\QQ[X]$.
\end{lemma}
\begin{proof}
  \leanok
  \uses{lem:evenpoly, lem:involution}
  Assume $f_n$ reducible, and let $m := \min\{k \mid f_k \text{ is reducible}\}$. Then
  $1 \le m \le n$, $f_{m-1}$ is irreducible, and $f_m = f_{m-1}(X^2 + a)$ is reducible. Let
  $\varepsilon := 1$ if $m > 1$ and $\varepsilon := -1$ if $m = 1$. Now $f_m$ is an even
  polynomial without nontrivial even divisors (since $f_{m-1}(Y + a)$ would be reducible).
  Therefore, if $h$ is a nontrivial divisor of $f_m$, then so is $h(-X)$, and $h$ and $h(-X)$ are
  coprime (since their g.c.d.\ is even). Thus, $f_m$ being reducible, there exists a polynomial
  $g$ with $\varepsilon \cdot f_m = g(X) \cdot g(-X)$. In particular,
  $c_m = \varepsilon f_m(0) = g(0)^2$ is a square in $\QQ$.
  \formalnote{The sign $\varepsilon$ is carried explicitly: what is produced is
  $(-1)^{\deg g} f_m = g \cdot g(-X)$ with $2 \deg g = \deg f_m$, and $(-1)^{\deg g}$ is exactly
  the sign relating $c_m$ to $f_m(0)$ (\lname{intCast\_cSeq\_eq\_ite\_mul\_eval\_zero}).
  ``No nontrivial even divisor'' is proved in the sharper form needed here: any divisor of
  $F \circ (X^2 + c)$ that is even \emph{up to associates} is a unit or an associate of the whole
  polynomial (\lname{isUnit\_or\_associated\_of\_dvd\_comp\_of\_associated}); the pairing of the
  irreducible factors under $X \mapsto -X$ is a fixed-point-free involution argument on the
  multiset of normalized factors (Lemma~\ref{lem:involution}).}
\end{proof}

\begin{corollary}
  \label{cor:1.3}
  \uses{lem:1.2}
  \lean{QuadraticIterates.irreducible_iteratedPoly_of_pos,
    QuadraticIterates.irreducible_iteratedPoly, QuadraticIterates.not_isSquare_cSeq,
    QuadraticIterates.not_isSquare_neg_of_irreducible}
  \leanok
  All $f_n$ are irreducible over $\QQ$.
\end{corollary}
\begin{proof}
  \leanok
  \uses{lem:1.1}
  The proof of Lemma~\ref{lem:1.1}~a) shows $|c_n| \ge |a|$ for all $n \ge 1$. Hence we have if
  $a > 0$: $c_n^2 < c_n^2 + a = c_{n+1} < (c_n + 1)^2$, so $c_{n+1}$ is not a square; if
  $a < -1$: $(c_n - 1)^2 < c_n^2 + a = c_{n+1} \le c_n^2$, so $c_{n+1}$ is not a square. Since
  $c_1 = -a$ is not a square either, Lemma~\ref{lem:1.2} applies.
  \formalnote{The formal statement covers $n = 0$ as well ($f_0 = X$ is irreducible). The
  converse implication is also formalized: irreducibility of a single $f_n$ with $n \ge 1$ forces
  $-a$ to be a non-square (\lname{not\_isSquare\_neg\_of\_irreducible}), since otherwise
  $f_n = (f_{n-1} - r)(f_{n-1} + r)$. This is what allows Lemma~\ref{lem:1.6} to be stated with
  irreducibility as its hypothesis.}
\end{proof}

The well-known example $a = -2$, where $K_n = \QQ(\zeta_{2^{n+2}}) \cap \mathbb{R}$, shows that
it is possible to have $\Om_n$ as small as $C_{2^n}$. We will see later that we can also give
examples where the other extreme, $\Om_n \cong \Ct{n}$, occurs.

\begin{lemma}
  \label{lem:1.4}
  \uses{def:galois, def:wreath, facts:1.0}
  \lean{QuadraticIterates.embed_equiv, QuadraticIterates.odoni_embedding,
    QuadraticIterates.isPGroup_galoisGroup, QuadraticIterates.pow_two_pow_apply_root,
    QuadraticIterates.ncard_rootSet_iteratedPoly_le,
    QuadraticIterates.nonempty_mulEquiv_of_finrank_eq}
  \leanok
  $\Om_n$ can be embedded into $\Ct{n}$. In particular, we have for all $n$:
  \[ \Om_{n+1} \cong \Ct{n+1} \iff \Om_n \cong \Ct{n} \text{ and } [K_{n+1} : K_n] = 2^{2^n}. \]
\end{lemma}
\begin{proof}
  \leanok
  \uses{lem:sylowwreath}
  Using the result of \cite{odoni}, p.~103, Cor.~1.2, and the fact that Galois groups do not
  increase under specialization, we get $\Om_n \hookrightarrow \Ct{n}$. The rest is then clear
  because of $\#\Ct{n+1}/\#\Ct{n} = 2^{2^n}$ and $[K_{n+1} : K_n] \le 2^{2^n}$
  (Facts~\ref{facts:1.0}).
  \formalnote{The embedding is proved directly rather than quoted, and without specializing from
  a generic polynomial. Every element of $\Om_n$ has order a power of $2$, because an
  automorphism raised to the power $2^m$ fixes each root of $f_m$ (induction on $m$, using that
  the roots of $f_{m+1}$ are the square roots of the shifted roots of $f_m$); so $\Om_n$ is a
  $2$-group acting faithfully on the at most $2^n$ roots of $f_n$, hence embeds into a Sylow
  $2$-subgroup of $S_{2^n}$, which is $\Ct{n}$ (Lemma~\ref{lem:sylowwreath}). The degree
  bookkeeping in the second part is done with $\#\Om_n = [K_n:\QQ]$ and
  $[K_{n+1}:\QQ] = [K_{n+1}:K_n]\,[K_n:\QQ]$.}
\end{proof}

\begin{lemma}
  \label{lem:1.5}
  \uses{def:twoindep, def:galois, def:cseq}
  \lean{QuadraticIterates.kummer_extension_criterion,
    QuadraticIterates.isSquare_algebraMap_iff_exists_mul_prod,
    QuadraticIterates.isSquare_algebraMap_iff_exists_sq_eq,
    QuadraticIterates.finrank_adjoin_range_eq_two_pow,
    QuadraticIterates.finrank_adjoin_le_two_pow,
    QuadraticIterates.card_monoidHom_eq_two_pow,
    QuadraticIterates.isSquare_algebraMap_cSeq,
    QuadraticIterates.twoIndependent_snoc_iff,
    QuadraticIterates.finrank_eq_of_nonempty_mulEquiv}
  \leanok
  If $\Om_n \cong \Ct{n}$ and $c_1, c_2, \dots, c_n$ are $2$-independent, then for all
  $c \in \QQ^*$:
  \[ c \notin (K_n^*)^2 \iff c_1, c_2, \dots, c_n, c \text{ are $2$-independent.} \]
\end{lemma}
\begin{proof}
  \leanok
  \uses{lem:wreathab, lem:multiquadratic, lem:descent}
  One proves easily by induction (using the fact that for a transitive permutation group $G$ one
  has $(G[H])^{ab} \cong G^{ab} \times H^{ab}$, where $G[H]$ denotes the wreath product) that
  $(\Ct{n})^{ab}$, the largest abelian factor group of $\Ct{n}$, is isomorphic with $C_2^n$.
  Thus, the largest $2$-Kummer extension of $\QQ$ within $K_n$ has degree $2^n$, and the kernel
  of the natural map $\QQ^*/(\QQ^*)^2 \to K_n^*/(K_n^*)^2$ has $\FF_2$-dimension $n$. Now,
  denoting by $A$ the set of zeros of $f_{m-1}$, we have that
  $c_m = f_{m-1}(-a) = \prod_{\alpha \in A} (-a - \alpha) = \prod_{\alpha \in A} (-a + \alpha)$
  is a square in $K_m$ (note that $A = -A$ and that for $\alpha \in A$, $\alpha - a$ is a square
  in $K_m$), so all of $c_1, \dots, c_n$ are squares in $K_n$. Since $c_1, \dots, c_n$ are
  $2$-independent, the above kernel must be generated by (the residue classes of)
  $c_1, \dots, c_n$. This implies the asserted equivalence.
  \formalnote{The two halves of ``the kernel has dimension exactly $n$'' are separate lemmas. For
  the lower bound, $\QQ(\sqrt{c_1}, \dots, \sqrt{c_n})$ has degree $2^n$ because the $c_i$ are
  $2$-independent (\lname{finrank\_adjoin\_range\_eq\_two\_pow}, via
  Lemma~\ref{lem:multiquadratic}). For the upper bound, any subfield of $K_n$ generated over
  $\QQ$ by square roots of rationals is Galois with elementary abelian group, so its degree is
  bounded by $\#\Hom(\Om_n, C_2) = 2^n$ (Lemma~\ref{lem:wreathab}). Hence a square root of $c$
  lying in $K_n$ already lies in $\QQ(\sqrt{c_1}, \dots, \sqrt{c_n})$, and the multiquadratic
  descent lemma (Lemma~\ref{lem:descent}) turns that into the statement that $c$ times a
  subproduct of the $c_i$ is a square in $\QQ$ --- which is exactly the failure of
  $2$-independence of $c_1, \dots, c_n, c$.}
\end{proof}

% The paper's remark after Lemma 1.5 is unnumbered, and it stays unnumbered here: a `remark`
% environment would consume a number and shift Lemma 1.6 to 1.7.
\emph{Remark.} We need to take $c_1 = -a$ instead of the more natural $c_1 = f(0) = a$ in order
to have $c_1$ a square in $K_1$ (\lname{cSeq\_one}).

\begin{lemma}[cf.\ \cite{odoni}, Lemma 4.3]
  \label{lem:1.6}
  \uses{def:galois, def:cseq, facts:1.0}
  \lean{QuadraticIterates.degree_criterion, QuadraticIterates.relfinrank_succ_eq_pow,
    QuadraticIterates.rootShift, QuadraticIterates.coe_rootShift,
    QuadraticIterates.rootShift_ne_zero, QuadraticIterates.prod_radicand_eq_cSeq,
    QuadraticIterates.prod_aroots_sub_eq_cSeq,
    QuadraticIterates.isPretransitive_galoisGroup,
    QuadraticIterates.exists_ringEquiv_radicand_smul,
    QuadraticIterates.degree_criterion_zero,
    QuadraticIterates.finrank_splittingField_one}
  \leanok
  For all $n$: $[K_{n+1} : K_n] = 2^{2^n} \iff c_{n+1}$ is not a square in $K_n$.
\end{lemma}
\begin{proof}
  \leanok
  \uses{cor:1.3, lem:multiquadratic, lem:invariant, lem:relations}
  Let $\alpha_j$ $(j = 1, 2, \dots, 2^n)$ be the roots of $f_n$ in $K_n$. Then
  $K_n \subset K_{n+1}$ is a $2$-Kummer extension, and $[K_{n+1} : K_n] = 2^{2^n - \dim V}$,
  where $V$ is the $\FF_2$-vector space of all $(\varepsilon_1, \dots, \varepsilon_{2^n}) \in
  \FF_2^{2^n}$ such that $\prod_j (\alpha_j - a)^{\varepsilon_j}$ is a square in $K_n$. We show
  that $V \ne 0 \iff (1, \dots, 1) \in V$: $\Om_n$ operates on $V$ by permuting the components of
  the elements of $V$ according to its action on the $\alpha_j$, thus $V$ becomes an
  $\FF_2[\Om_n]$-module. Now, if $G$ is a $2$-group and $M \ne 0$ is an $\FF_2[G]$-module, then
  the submodule $M^G$ of $G$-invariant elements of $M$ is also nontrivial. (Induction on $\#G$:
  Let $G \cong C_2 = \{1, \sigma\}$ and $0 \ne y \in M$. Either $y = \sigma y$, then
  $0 \ne y \in M^G$, or $y \ne \sigma y$, then $0 \ne y + \sigma y \in M^G$. If $\#G > 2$, take a
  nontrivial normal subgroup $N$ of $G$. Then $M$ is an $\FF_2[N]$-module, too, so by induction
  $M^N \ne 0$. Now $M^N$ is an $\FF_2[G/N]$-module, so $M^G = (M^N)^{G/N} \ne 0$.) So, if
  $V \ne 0$, then $V^{\Om_n} \ne 0$. Since $\Om_n$ operates transitively on the $\alpha_j$
  ($f_n$ is irreducible by Cor.~\ref{cor:1.3}), the only nontrivial invariant element in $V$ is
  $(1, \dots, 1)$.

  From the above it is clear that the left hand side is equivalent to
  $(1, \dots, 1) \notin V$, whereas $(1, \dots, 1) \in V$ means that
  $\prod_j (\alpha_j - a) = c_{n+1}$ is a square in $K_n$.
  \formalnote{The hypothesis used is that $f_n$ is irreducible, which by
  Corollary~\ref{cor:1.3} holds under the standing assumption on $a$; the paper instead argues
  under the maximality hypothesis $[K_n : \QQ] = 2^{2^n - 1}$, which is strictly stronger. The
  degenerate case $n = 0$ is included ($K_1 = \QQ(\sqrt{-a})$ has degree $2$ over $\QQ$ exactly
  when $-a$ is not a square). $V$ is \texttt{rootRelations} of the family
  $\beta \mapsto \beta - a$ of shifted roots (\texttt{rootShift}); that this family consists of
  nonzero elements is where $c_{n+1} \neq 0$ enters. The identity
  $\prod_j (\alpha_j - a) = c_{n+1}$ is proved by evaluating the monic split $f_n$ at $a$
  (\lname{prod\_aroots\_sub\_eq\_cSeq}).}
\end{proof}

Summing up, we can now prove the following

\begin{theorem}[Section 1]
  \label{thm:section1}
  \uses{def:galois, def:wreath, def:twoindep, def:cseq, def:bseq}
  \lean{QuadraticIterates.section1_equiv, QuadraticIterates.section1_a_iff_b,
    QuadraticIterates.section1_b_iff_c, QuadraticIterates.section1_squarefree,
    QuadraticIterates.sqClass_cSeq_eq_sum_divisors,
    QuadraticIterates.sqClass_bSeq_eq_sum_divisorsAntidiagonal}
  \leanok
  \begin{enumerate}
    \item For all $n$, the following statements are equivalent:
      \begin{enumerate}
        \item[a)] $\Om_n \cong \Ct{n}$;
        \item[b)] $c_1, c_2, \dots, c_n$ are $2$-independent;
        \item[c)] $b_1, b_2, \dots, b_n$ are $2$-independent.
      \end{enumerate}
    \item If none of $|b_2|, |b_3|, \dots, |b_n|$ is a square in $\QQ$, then $\Om_n \cong \Ct{n}$.
  \end{enumerate}
\end{theorem}
\begin{proof}
  \leanok
  \uses{lem:1.1, lem:1.4, lem:1.5, lem:1.6, cor:1.3, lem:coprimesquare}
  1) By induction on $n$. $n = 0$ is trivial. So let $n \ge 1$ and suppose 1) holds for $n - 1$.

  ``a) $\Leftrightarrow$ b)'': $\Om_n \cong \Ct{n}$ $\iff$ (by \ref{lem:1.4})
  $\Om_{n-1} \cong \Ct{n-1}$ and $[K_n : K_{n-1}] = 2^{2^{n-1}}$ $\iff$ (by induction hypothesis
  and \ref{lem:1.6}) $c_1, \dots, c_{n-1}$ are $2$-independent and $c_n \notin (K_{n-1}^*)^2$
  $\iff$ (by \ref{lem:1.5} and induction hypothesis) $c_1, \dots, c_n$ are $2$-independent.

  ``b) $\Leftrightarrow$ c)'': This is clear because of $c_m = \prod_{d \dvd m} b_d$.

  2) Since by Lemma~\ref{lem:1.1} $b_1, b_2, \dots, b_n$ are pairwise coprime, 1)~c) is
  equivalent to: none of the $b_m$ is a square, and at most one of the $-b_m$ is a square. Since
  $b_1 = -a$ is not a square, condition 2) implies 1)~c).
  \formalnote{In 1), ``b) $\Leftrightarrow$ c)'' is run in $\QQ^*/(\QQ^*)^2$: the classes of the
  $c_m$ and of the $b_m$ span the same $\FF_2$-subspace (each is a $\ZZ$-combination of the
  others, by Möbius inversion in both directions), and $2$-independence of a family of $n$
  vectors is equivalent to that span having dimension $n$. In 2), the passage from a square
  subproduct to a single square factor is the pairwise coprimality of the $b_m$: a square product
  of pairwise coprime integers has all $|b_m|$ square (Lemma~\ref{lem:coprimesquare}); the index
  $m = 1$ is excluded separately using $b_1 = -a$.}
\end{proof}
